% !TEX root = ../tcc.tex
% Resumo: ate ~250 palavras, sem citacoes nem notas de rodape (regra IEEE).

\begin{abstract}
Em pipelines de IA generativa com vários estágios, a saída de um estágio vira a entrada do próximo, e um erro introduzido cedo pode chegar à resposta final sem que se saiba onde começou.
Este trabalho investiga como o contrato de handoff, o formato da informação trocada entre estágios, afeta a propagação, a correção e a atribuição de origem desse erro.
Um pipeline fixo de três estágios (Extrator, Analista e Redator) roda sobre o FinanceBench com quatro contratos: texto livre (A), tipado (B), tipado com lineage (C) e tipado com lineage e o trecho de evidência da fonte (D).
Um injetor determinístico corrompe um claim por execução, e o mesmo estado corrompido é projetado nas quatro condições.
O protocolo foi repetido com três modelos: GPT-5.6 Luna, Claude Haiku 4.5 e o SLM Qwen3.5 9B.
Nos três, a estrutura sozinha não melhorou a acurácia final e aumentou a profundidade de propagação.
O lineage elevou a atribuição correta do claim de origem de menos de 0,09 para cerca de 0,4, mas não reduziu a sobrevivência do erro.
Nos dois modelos de API, só a evidência reduziu a sobrevivência e a propagação do erro, e a atribuição chegou a cerca de 0,8 (+0,38 a +0,43 sobre C), com quase quatro vezes mais tokens que o texto livre.
No Qwen, a evidência melhorou a acurácia final, mas não a contenção nem a atribuição.
Estrutura e lineage, sozinhos, não tornam o pipeline mais confiável: o ganho vem da evidência que permite conferir a proveniência, e depende de o modelo saber usá-la.
\end{abstract}

\begin{IEEEkeywords}
pipelines de IA generativa, propagação de erros, sistemas multiagente, observabilidade, atribuição de causa raiz, contratos de handoff
\end{IEEEkeywords}
